\documentclass[a4paper,12pt]{article}
\usepackage[hmargin=2cm,top=4cm,headheight=65pt,footskip=45pt]{geometry}
\usepackage[utf8]{inputenc}
\usepackage{graphicx}
\usepackage[hidelinks]{hyperref}
\usepackage{array}
\usepackage{lastpage}
\usepackage{lipsum}
\usepackage{fancyvrb}
\usepackage{color}
\usepackage{fancyhdr}
\usepackage{amsmath}
\usepackage{enumitem}
\usepackage{titlesec}

\definecolor{customGray}{RGB}{128,128,128}
%\definecolor{Eblue}{RGB}{0,21,87}
%\definecolor{Eblue}{rgb}{0.2, 0.2, 0.6}
\definecolor{Eblue}{rgb}{0.0, 0.18, 0.39}
%==============Header & Footnote==============

\pagestyle{fancy}
\renewcommand{\headrulewidth}{0pt}
\fancyhead[C,CO,L,LO,R,RO]{}
\fancyhead[C]{%
          \begin{tabular}{|m{3.0cm}|m{10.0cm}|m{2.5cm}|}
          \hline
          \centering\vspace{1.75mm}\includegraphics[scale=0.275]{logo.pdf} &
          \centering
          {\footnotesize {\sf UNIVERSIDAD EAFIT\\ SCHOOL OF ENGINEERING\\
          \vspace{-1mm}DEPARTMENT OF SYSTEMS AND INFORMATICS}} &
          \centering
          \footnotesize{Page \thepage\ de \pageref{LastPage}\\
          ST245\\
          \vspace{-0.75mm}Data Structures
          }\tabularnewline
          \hline
          \end{tabular}%
}
\fancyfoot[C,CO,L,LO,R,RO]{}
\fancyfoot[C]{
          \begin{centering}
            \textcolor{customGray}{{\footnotesize {\sf Profoessor Mauricio Toro Bermudez\\
            Phone: $(+57) (4) 261 95 00$ Ext. $9473$. Office: $19 - 627$\\
            \vspace{-1mm}E-mail: mtorobe@eafit.edu.co}}}
        \end{centering}
}

%=============CustomEnumItem===========

\setlist[enumerate]{label=\color{Eblue}\textbf{\roman*.}}

%=============CustomSecSubSec==========

\titleformat{\section}[hang]
{\normalsize\bfseries\itshape\color{black}}{\bfseries\itshape\color{Eblue}\thesection)}{2.5mm}{}

\titleformat{\subsection}[hang]
{\normalsize\bfseries\itshape\color{black}}{\bfseries\color{Eblue}\thesection.\alph{subsection}.}{2.5mm}{}

%==============Title==============

\title{\color{Eblue}\textbf{Laboratory practice No. 1: Recursion}}
\author{
  \textbf{Juan S. Cárdenas Rodríguez}\\
  Universidad EAFIT\\
  Medellín, Colombia\\
  jscardenar@eafit.edu.co
\and
  \textbf{David Plazas Escudero}\\
  Universidad EAFIT\\
  Medellín, Colombia\\
  dplazas@eafit.edu.co
}

%=============Document=============
\begin{document}
  \maketitle
  \thispagestyle{fancy}

  \section{ONLINE EXERCISES (CODINGBAT)}
  \subsection{Recursion I}
    \begin{enumerate}
      \item \begin{Verbatim}
      public int countPairs(String str) {             // c0
        if (str.length() <= 2) {                      // c1
          return 0;                                   // c2
        } else if (str.charAt(0) == str.charAt(2)) {  // c3
          return 1 + countPairs(str.substring(1));    // c4 + T(n-1)
        } else {                                      // c4
          return countPairs(str.substring(1));        // T(n-1)
        }
      }
      \end{Verbatim}
      Complexity of \texttt{countPairs} can be writen as:
      \begin{equation*}
        T\left( n \right)=\left\{\begin{array}{cc} c_{0}+c_{1}+c_{2} & n\leq 2 \\ c_{3}+c_{4}+T\left( n-1 \right) & n>2\end{array}\right.
      \end{equation*}
      Solving the recursive equation yields:
      \begin{equation*}
        T\left( n \right)=(c_3+c_4)n+k
      \end{equation*}
      Therefore, $T(n)$ is $O(cn+k)$ and applying the sum and product rule $T(n)$ is $O(n)$.
      \item \begin{Verbatim}
      public int countHi2(String str) {               // c0
        if (str.length() == 1 || str.length() == 0) { // c1
          return 0;                                   // c2
        } else if (str.charAt(0) == 'x') {            // c3
          if (str.charAt(1) == 'h'
          && str.charAt(2) == 'i') {                  // c4
            return countHi2(str.substring(2));        // T(n-2)
          } else {                                    // c5
            return countHi2(str.substring(1));        // T(n-1)
          }
        } else if (str.charAt(0) == 'h'
          && str.charAt(1) == 'i') {                  // c5
          return 1 + countHi2(str.substring(1));      // c5
        } else {                                      // c6
          return countHi2(str.substring(1));          // T(n-1)
        }
      }
      \end{Verbatim}
      The complexity of \texttt{countHi2} can be writen as:
      \begin{equation*}
        T\left( n \right)=\left\{\begin{array}{cc} c_{0}+c_{1}+c_{2} & n\leq 1 \\ c_5+T\left( n-1 \right) & n>1\end{array}\right.
      \end{equation*}
      Solving the recursive equation for this algorithm, yields:
      \begin{equation*}
        T\left( n \right)=c_5n+k
      \end{equation*}
      Then, $T(n)$ is $O(c_5n+k)$ and applying the sum and product rule $T(n)$ is $O(n)$.
      \item \begin{Verbatim}
      public int countAbc(String str) {               // c0
        if (str.length() == 0 || str.length() == 1
        || str.length() == 2) {                       // c1
          return 0;                                   // c2
        } else if (str.charAt(0) == 'a'
          && str.charAt(1) == 'b'
          && (str.charAt(2) == 'c'
          || str.charAt(2) == 'a')) {                 // c3
          return 1 + countAbc(str.substring(1));      // c4 + T(n-1)
        } else {                                      // c5
          return countAbc(str.substring(1));          // T(n-1)
        }
      }
      \end{Verbatim}
      The complexity of \texttt{countAbc} can be writen as:
      \begin{equation*}
        T\left( n \right)=\left\{\begin{array}{cc} c_{0}+c_{1}+c_{2} & n\leq 2 \\ c_3 + c_ 4 + T\left( n-1 \right) & n>2\end{array}1\right.
      \end{equation*}
      The solution to this recursive equation yields:
      \begin{equation*}
        T(n)=(c_3+c_4)n+k
      \end{equation*}
      Therefore, $T(n)$ is $O((c_3+c_4)n+k)$ and applying the sum and product rule $T(n)$ is $O(n)$.

      \item \begin{Verbatim}
      public String parenBit(String str) {            // c0
        int a = str.length();                         // c1
        if (a <= 1) {                                 // c2
            return "";                                // c3
        }
        if (str.substring(a - 1).equals(")")) {       // c4
          int paren = str.indexOf("(");               // c5
          return str.substring(paren);                // T(n-k)
        }
        return parenBit(str.substring(0,a - 1));      // T(n-1)
      }
      \end{Verbatim}
      The complexity of \texttt{parenBit} can be writen as:
      \begin{equation*}
        T\left( n \right)=\left\{\begin{array}{cc} c_{0}+c_{1}+c_{2}+c_{3} & n\leq 1 \\ c_{4}+T\left( n-1 \right) & n>1\end{array}\right.
      \end{equation*}
      Solving the recursive equation yields:
      \begin{equation*}
        T(n)=c_4n+k
      \end{equation*}
      Then $T(n)$ is $O(c_4n+k)$ and applying the product and sum rules, we obtain that $T(n)$ is $O(n)$.
      \item \begin{Verbatim}
      public int strCount(String str, String sub) {
        int a = str.length();
        int b = sub.length();
        if (a < b || b == 0){
          return 0;
        }
        if (str.substring(a - b).equals(sub)) {
          return 1 + strCount(str.substring(0,a - b),sub);
        }
        return strCount(str.substring(0,a - 1), sub);
      }
      \end{Verbatim}
    \end{enumerate}

    \subsection{Recursion II}
    \begin{enumerate}
      \item \begin{Verbatim}
      public boolean splitArray(int[] nums) {
        return splitArrayAux(nums, 0, 0, 0);
      }
      public boolean splitArrayAux(int [] nums, int start,
        int first, int second) {                    // c1
        if (start == nums.length) {                 // c2
          return first == second;                   // c3
        } else {                                    // c4
          return splitArrayAux(nums, start + 1,
            first + nums[start], second)
          || splitArrayAux(nums, start + 1, first,
            second + nums[start]);                  // c5 + 2T(n-1)
        }
      }
      \end{Verbatim}
      Complexity of \texttt{splitArray} can be writen as:
      \begin{equation*}
        T\left(n\right)=\left\{\begin{array}{cc}c_1+c_2+c_3+c_4&n=0\\c_5+2T\left(n-1\right)&n\neq0\end{array}\right.
      \end{equation*}
      The solution to this recursive equation is:
      \begin{equation*}
        T\left( n \right)=k2^{n-1} + \left( 2^{n} - 1 \right) \left( c1 + c2 + c3 + c4 + c5 \right)
      \end{equation*}
      Then, $T\left(n\right)$ is $O\left(k2^{n-1} + \left( 2^{n} - 1 \right) \left( c_1 + c_2 + c_3 + c_4 + c_5 \right)\right)$.\\
      Therefore, applying the sum and product rule $T(n)$ is $O(2^n)$.

      \item \begin{Verbatim}
      public boolean splitOdd10(int[] nums) {
        return splitOdd10Aux(nums, 0, 0, 0);
      }
      public boolean splitOdd10Aux(int [] nums, int start,
        int first, int second) {                          // c1
        if (start == nums.length) {                       // c2
          return (first % 10 == 0) && (second % 2 != 0);  // c3
        } else {                                          // c4
          return splitOdd10Aux(nums, start + 1;
            first + nums[start], second) ||
          splitOdd10Aux(nums, start + 1,
            first, second + nums[start]);                 // c5 + 2T(n-1)
        }
      }
      \end{Verbatim}
      Complexity of \texttt{splitOdd10} can be writen as:
      \begin{equation*}
        T\left(n\right)=\left\{\begin{array}{cc}c_1+c_2+c_3+c_4&n=start\\c_5+2T\left(n-1\right)&n\neq start\end{array}\right.
      \end{equation*}
      The solution to this recursive equation is:
      \begin{equation*}
        T\left( n \right)=k2^{n-1} + \left( 2^{n} - 1 \right) \left( c1 + c2 + c3 + c4 + c5 \right)
      \end{equation*}
      Then, $T\left(n\right)$ is $O\left(k2^{n-1} + \left( 2^{n} - 1 \right) \left( c1 + c2 + c3 + c4 + c5 \right)\right)$.\\
      Therefore, applying the sum and product rule $T(n)$ is $O(2^n)$.
      \item \begin{Verbatim}
      public boolean groupSumClump(int start, int[] nums,
        int target) {                                    // c1
        if (start >= nums.length) {                      // c2
          return target == 0;                            // c3
        }
        int sum = 0;                                     // c4
        int i;                                           // c5
        for (i = start; i < nums.length; i++) {          // c6 * n
          if (nums[i] == nums[start]) {                  // c7 * n
            sum += nums[start];                          // c8 * n
          } else {                                       // c9 * n
            break;                                       // c10 * n
          }
        }
        return groupSumClump(i, nums, target - sum)
        || groupSumClump(i, nums, target);               // c12 + 2T(n-1)
      }
      \end{Verbatim}
      Can be writen as:
      \begin{equation*}
        T\left(n\right)=\left\{\begin{array}{cc}c_{3}&n\leq \text{start}\\c_{1}+c_{2}+c_{4}+c_{5}+c_{12}+\left(c_{6}+c_{7}+c_{8}+c_{9}+c_{10}\right)n+2T\left(n-1\right) & n>start\end{array}\right.
      \end{equation*}
      The solution to this recursive equation is:
      \begin{equation*}
        T\left( n \right) = 2^{n-1} \left( c + 4 c_1 \right) + c_2 \left( 2^{n} - 1 \right) - c_1 \left( n +2 \right).
      \end{equation*}
      Then, $T\left(n\right)$ is $O\left(2^{n-1} \left( c + 4 c_1 \right) + c_2 \left( 2^{n} - 1 \right) - c_1 \left( n +2 \right)\right)$\\
      Therefore, applying the sum and product rule $T(n)$ is $O(2^n)$.

      \item \begin{Verbatim}
      public boolean groupSum5(int start, int[] nums, int target) {
        if (start == nums.length) {                     // c1
          return target == 0;                           // c2
        } else {                                        // c3
          if (nums[start] % 5 == 0) {                   // c4
            return groupSum5(start + 1, nums,
            target - nums[start]);                      // c5 + T(n-1)
          } else if (start > 0 && nums[start] == 1
            && nums[start - 1] % 5 == 0) {              // c6
            return groupSum5(start + 1, nums, target);  // c7 + T(n-1)
          } else {                                      // c8
            return groupSum5(start + 1, nums,
            target - nums[start])
            || groupSum5(start + 1, nums, target);      // c9 + 2T(n-1)
          }
        }
      }
      \end{Verbatim}
      Taking into account that the case $c_9+2T(n-1)$ is the worst out of all, we can write the recursive equation as:
      \begin{equation*}
        T\left( n \right)=\left\{\begin{array}{cc} c_{2} & n=start \\ c_{1}+c_{3}+c_{4}+c_{6}+c_{8}+c_{9}+2T\left( n-1 \right) & n\neq start\end{array}\right.
      \end{equation*}
      The solution to this recursive equation is:
        \begin{equation*}
T\left( n \right)=\left( c_{1}+c_{3}+c_{4}+c_{6}+c_{8}+c_{9} \right)\left( 2^{n}-1 \right)+c2^{n-1}
        \end{equation*}
      \item \begin{Verbatim}
      public boolean split53(int[] nums) {
        return split53Aux(nums, 0, 0, 0);
      }
      public boolean split53Aux(int [] nums, int start,
        int first, int second) {
        if (start == nums.length) {
          return first == second;
        } else {
          if (nums[start] % 5 == 0) {
            return split53Aux(nums, start + 1, first + nums[start], second);
          } else if (nums[start] % 3 == 0) {
            return split53Aux(nums, start + 1, first, second + nums[start]);
          } else {
            return split53Aux(nums, start + 1, first + nums[start], second)
            || split53Aux(nums, start + 1, first, second + nums[start]);
          }
        }
      }
      \end{Verbatim}
    \end{enumerate}\cite{Wolfram}

    \section{ArrayMax}
    



    \section{What did you learn about \texttt{Stack Overflow}?}
      The Stack Overflow error is caused by a bad recursive call -for example you do not make
      the problem simpler every time you make a recursive call- or when you do not have a stopping
      condition.\cite{StackOverflow} Java Stack memory is used for execution of a thread.
      Whenever a method is invoked, a new block is created in the stack memory for the method to hold
      local primitive values and reference to other objects in the method.\cite{HeapStack}


    \newpage
    \bibliographystyle{plain}
    \bibliography{Lab.bib}
\end{document}
